
\textbf{Proposition 2.3.1.} \uline{Pour tout schéma propre} $X$ \uline{sur} $k$\uline{, et pour tout} $y\in\Filt_j(X)$\uline{, l'application}
\[
S_y:\Pic(X)\longrightarrow\mathbb Z\hookrightarrow\mathbb Q
\]
\uline{définie par}
\[
\cl_P(\underline{L})\rightsquigarrow g_*(\cl^{\bullet}(\underline{L})y)=\chi(\cl^{\bullet}(\underline{L})y
\]
\uline{est une application polynôme de degré} $\leq j$.

Lorsque $\underline{L}_1,\ldots,\underline{L}_n$ sont des $\underline{O}_X$-modules inversibles et $\underline{F}$ un $\underline{O}_X$-module cohérent tel que $\dim(\Supp(\underline{F}))=j$, alors 2.3.1. implique en particulier que la caractéristique d'Euler-Poincaré
\begin{equation}\tag{2.3.2}
\chi(\underline{L}_1^{\otimes m_1}\otimes\cdots\otimes \underline{L}_n^{\otimes m_n}\otimes \underline{F})
\end{equation}
est un polynôme numérique dans les variables $m_1,\ldots,m_n\in\mathbb Z$, de degré $\leq\dimdot\operatorname{supp.}F$, parfois appellé \emph{polynôme de Snapper} pour $\underline{F}$ et la suite $\underline{L}_1,\ldots,\underline{L}_n$.

\subsection*{2.4.}
Rappelons [2] que le foncteur de Picard $\text{\uwave{Pic}}_{X/k}$ de $X$ sur $k$ est le faisceau en groupes dans la topologie ($\fpqc$) de $k$ associé au préfaisceau
\[
T\rightsquigarrow \Pic(X\times_kT)\qquad (T\in\operatorname{Ob}(\Sch/k)).
\]
Lorsque $k$ est algébriquement clos, on sait que l'homomorphisme canonique
\[
\Pic(X)\longrightarrow\Pic_{X/k}(\Spec(k))
\]
est un isomorphisme.



% Local continuation macros
\providecommand{\Filt}{\operatorname{Fil}}
\providecommand{\fpqc}{\mathrm{fpqc}}
\providecommand{\Sch}{\operatorname{Sch}}
\renewcommand{\Ob}{\operatorname{ob}}
\providecommand{\Gm}{\mathbb G_m}
\providecommand{\Der}{\operatorname{Der}}
\providecommand{\SL}{\operatorname{SL}}
\providecommand{\Aut}{\operatorname{Aut}}
\providecommand{\Mod}{\operatorname{Mod}}
\providecommand{\Tors}{\operatorname{Tors}}
\providecommand{\Mat}{\operatorname{Mat}}
\providecommand{\len}{\operatorname{long}}
\renewcommand{\id}{\operatorname{id}}
\providecommand{\Spec}{\operatorname{Spec}}
\providecommand{\Picf}{\operatorname{Pic}}
\providecommand{\Ens}{\operatorname{Ens}}
\providecommand{\clp}{\operatorname{cl}_{P}}
\providecommand{\clo}{\operatorname{cl}^{\circ}}
\providecommand{\clb}{\operatorname{cl}_{\bullet}}
\providecommand{\Grb}{\operatorname{Gr}_{\bullet}}
\providecommand{\Gro}{\operatorname{Gr}^{\bullet}}
\providecommand{\longO}[1]
